\documentclass[11pt]{article}

% Standard, Overleaf-safe packages only
\usepackage[utf8]{inputenc}
\usepackage[T1]{fontenc}
\usepackage[margin=1in]{geometry}
\usepackage{amsmath,amssymb}
\usepackage{graphicx}
\usepackage{booktabs}
\usepackage{xcolor}
\usepackage{caption}
\usepackage{subcaption}
\usepackage[numbers,sort&compress]{natbib}
\usepackage[hidelinks]{hyperref}

\graphicspath{{figures/}{../results/figures/}}

% Convenience macros
\newcommand{\libgoal}{LIBERO-Goal}
\newcommand{\libten}{\texttt{libero\_10}}
\newcommand{\psft}{\pi_{\mathrm{SFT}}}
\newcommand{\pcsucc}{\mathrm{pc\_success}}

\title{\textbf{Reward-Weighted Behavior Cloning for Flow-Matching\\
Vision-Language-Action Policies: A Dual-Suite Null on LIBERO}}

\author{
  Gaurav Mitra\\
  Georgia Institute of Technology\thanks{Correspondence: \texttt{mitragaurav10@gmail.com}. Code and
  all result files: \url{https://github.com/Gogo2015/VLATune}.}
}
\date{}

\begin{document}
\maketitle

\begin{abstract}
Supervised fine-tuning (SFT) of vision-language-action (VLA) models tends to plateau, and
reinforcement learning (RL) is the obvious next step. Modern VLAs, though, emit actions through a
flow-matching denoiser with no tractable per-action log-probability, so PPO does not apply. We
study a lightweight alternative that fits flow matching: filtered, reward-weighted behavior
cloning on a policy's own successful rollouts, regularized toward the SFT policy by a vector-field
distillation anchor that stands in for an intractable marginal KL. Using SmolVLA (450M) on LIBERO
and a compute-frugal pipeline (a consumer GPU plus preemptible cloud, about \$20--60 total), we
evaluate the method in two regimes under a paired \emph{collection-equals-evaluation} protocol:
the saturated, in-distribution \libgoal{} suite (about 89\%) and an unsaturated, demo-seeded
\libten{} (LIBERO-Long, about 44\%). In neither regime does the method improve the policy durably.
Success rises for one iteration and returns to baseline the next ($89.0 \to 91.0 \to 89.4$ on
\libgoal{}; $44.0 \to 47.5 \to 42.5$ on \libten{}), a within-noise redistribution rather than a
learned gain. Two mechanisms explain the pattern: a marginal-task rescue-then-reversal, and a
starvation floor for tasks with zero seed successes. The null is informative because it holds even
with real headroom and no filter starvation. The contribution is the method, the paired protocol,
a map of where lightweight RL can and cannot move a flow-matching VLA, and the dual-suite null
itself.
\end{abstract}

\section{Introduction}

Generalist robot manipulation increasingly relies on \emph{vision-language-action} (VLA) models: a
single policy mapping camera images and a natural-language instruction to motor actions across
many tasks~\citep{kim2024openvla,shukor2025smolvla}. Two facts shape the research loop. First, the
strongest VLAs are large (OpenVLA is 7B~\citep{kim2024openvla}), which puts them out of reach of
modest compute. Second, improving a VLA with reinforcement learning is unsolved for the modern
variants whose action head is a flow-matching or diffusion model~\citep{chi2023diffusion,lipman2023flow}:
such a head produces actions by iteratively denoising noise, so the marginal likelihood of a
sampled action chunk is intractable and PPO's probability ratio~\citep{schulman2017ppo} cannot be
formed.

SmolVLA~\citep{shukor2025smolvla} addresses the first fact. At 450M parameters (about 100M of them
a trainable action expert, with a frozen VLM backbone) it is competitive on LIBERO while training
on a single consumer-class GPU. That makes the whole SFT-to-RL loop affordable, and turns the
second fact into a concrete, runnable question:

\begin{quote}
\emph{Can a lightweight, flow-matching-compatible RL step (filtered, reward-weighted behavior
cloning on the policy's own successes, anchored to the SFT policy by a vector-field trust-region)
improve a SmolVLA VLA beyond supervised fine-tuning on LIBERO?}
\end{quote}

We pick an objective that never needs an action log-probability: reward-weighted
regression~\citep{peters2007rwr,peng2019awr}, adapted to flow matching~\citep{orwcfm2025}. This
sidesteps the intractable-likelihood problem instead of working around it, as the noise-injection
and two-layer-MDP methods do~\citep{ren2024dppo,chen2025pirl,reinflow2025,flowgrpo2025}. We then
test the method as carefully as the compute budget allows and report what we find.

\paragraph{Finding.} Across both a saturated in-distribution suite and an unsaturated but
bootstrappable one, the method does not improve the policy durably. Aggregate success rises in the
first RL iteration, and the second erases it; the per-task gains reverse on the same tasks that
drove them. The up-then-down trajectory is itself the evidence that the bump is noise, not signal.

\paragraph{Contributions.}
\begin{enumerate}\itemsep2pt
  \item A tractable, flow-matching-native RL recipe (reward-weighted filtered BC with a
  vector-field distillation anchor as a KL surrogate), built as a custom fine-tune step on top of
  SmolVLA and validated by a gated smoke ladder (Section~\ref{sec:method}).
  \item A paired \emph{collection-equals-evaluation} protocol that turns each iteration's rollout
  collection into that policy's per-task evaluation at no extra cost, with identical initial states
  across iterations (Section~\ref{sec:method}).
  \item A dual-suite null: no durable gain on saturated \libgoal{} or on demo-seeded \libten{}.
  This is stronger than a single-suite null, because it holds with real headroom and no filter
  starvation (Section~\ref{sec:results}).
  \item Two mechanisms behind the null, rescue-then-reversal of marginal tasks and a starvation
  floor, and a map of where the method can and cannot move the policy (Section~\ref{sec:discussion}).
\end{enumerate}

\section{Background and Related Work}

\paragraph{SmolVLA and flow-matching action experts.}
SmolVLA~\citep{shukor2025smolvla} pairs a compact vision-language backbone with a small (about
100M) transformer action expert that produces an action \emph{chunk} via flow
matching~\citep{lipman2023flow}: starting from Gaussian noise, it integrates a learned velocity
field to a clean action. Training regresses the predicted velocity onto the noise-to-action
target, a supervised MSE with no normalizing constant and no marginal action likelihood. During
fine-tuning the backbone is frozen and only the action expert is trained, which is what makes the
model data- and compute-efficient.

\paragraph{LIBERO.}
LIBERO~\citep{liu2023libero} is a simulated (MuJoCo) manipulation benchmark with sub-suites
including Goal (10 tasks that share a scene but vary the instruction, isolating goal understanding)
and Long (\libten{}, 10 long-horizon multi-step tasks). Success is outcome-only and binary: the
simulator checks a ground-truth goal predicate at the step limit. The paper reports 92\% for
SmolVLA on \libgoal{}, while independent reproductions in the public LeRobot
codebase~\citep{lerobot2024} land near 78--84\%~\citep{nanovla2025}. We plan around that gap rather
than chase it.

\paragraph{Why not PPO; reward-weighted regression.}
PPO~\citep{schulman2017ppo} needs the ratio $\pi_\theta(a\mid s)/\pi_{\text{old}}(a\mid s)$, and so
a tractable action likelihood, which a flow-matching policy does not provide cheaply. Recent work
recovers a likelihood to run on-policy policy gradient, either by injecting noise for a per-step
Gaussian log-prob (ReinFlow~\citep{reinflow2025}, $\pi$RL~\citep{chen2025pirl},
Flow-GRPO~\citep{flowgrpo2025}) or by treating the denoising chain as an inner MDP
(DPPO~\citep{ren2024dppo}). We take the other route and choose an objective that never needs one:
reward-weighted, or advantage-weighted, regression~\citep{peters2007rwr,peng2019awr}, whose
flow-matching instantiation is online reward-weighted fine-tuning of flow
matching~\citep{orwcfm2025}. RL does move flow and VLA policies in other settings, often on
generalization~\citep{simplevlarl2025,rftflow2025,chen2025pirl}; our question is whether the
lightweight, likelihood-free version moves \emph{this} policy on LIBERO.

\section{Method}
\label{sec:method}

\subsection{Supervised fine-tuning (the RL launch point)}
We fine-tune \texttt{lerobot/smolvla\_base} on the \libgoal{} episodes of the
\texttt{HuggingFaceVLA/libero} dataset (428 usable episodes, batch 64, with the VLM backbone frozen
by default). Five non-obvious fixes were needed, and we document them for reproducibility: (i)
disabling the Hub push that otherwise aborts training; (ii) working around a broken episode-to-file
index in the dataset, which forces a full download and then an in-row \texttt{episode\_index}
filter; (iii) a camera-key contract (\texttt{image}/\texttt{image2} to
\texttt{camera1}/\texttt{camera2}) via \texttt{rename\_map}; (iv) the same \texttt{rename\_map} at
evaluation; and (v) pre-creating the LIBERO config to avoid an \texttt{EOFError} in headless
subprocesses. The run was preempted at 16k of 20k steps with converged loss (about 0.10); we
evaluate the 16k checkpoint and pin it as the immutable RL reference $\psft$.

\subsection{The RL loop: rollout, filter, reward-weighted fine-tune}
One outer iteration, starting from the current policy $\pi_k$ (with $\pi_0=\psft$):
\begin{enumerate}\itemsep2pt
  \item \textbf{Collect.} Run $N$ rollouts per task in LIBERO under the exact eval contract
  (relative control, \texttt{n\_action\_steps}=10, \texttt{rename\_map}, \texttt{MUJOCO\_GL=egl}).
  Record each step's observation and executed action chunk, the terminal binary success
  $r\in\{0,1\}$, and the step count.
  \item \textbf{Filter and weight.} Keep $(\text{obs},\text{action})$ pairs from successful
  trajectories. Weight each by a time-to-success factor $w=\exp(-\,\text{steps}/\tau_t)$, where the
  per-task $\tau_t$ is the median success length, normalized globally to mean 1.
  \item \textbf{Fine-tune.} Continue-train $\pi_k$ on the kept set with the flow-matching loss
  scaled per-sample by $w$, plus the anchor below, at a small learning rate ($2\times 10^{-5}$) for
  1000 steps, giving $\pi_{k+1}$.
  \item \textbf{Gate.} Evaluate $\pi_{k+1}$ and stop early on regression.
\end{enumerate}

\paragraph{The sparse-binary collapse.}
With $r\in\{0,1\}$ and no other signal, the return is binary, so reward-weighting is
mathematically identical to \emph{filtering}: keep successes, drop failures, behavior-clone the
rest. Under sparse binary reward our ``reward-weighted BC'' is rejection-sampling fine-tuning, the
robotics analogue of best-of-$N$ followed by SFT on the passes. The time-to-success weight $w$ is
the one cheap, non-degenerate reward shape that makes the weighting more than a filter, and we do
not overstate it. What makes this RL rather than more SFT is the signal source: the environment's
reward on the policy's own rollouts, not human demonstrations.

\paragraph{The vector-field anchor (a tractable KL surrogate).}
The exact marginal $\mathrm{KL}(\pi_\theta\,\Vert\,\psft)$ over action chunks is intractable for the
same reason PPO is. We anchor the mechanism instead. Keeping a frozen copy of $\psft$, we penalize
the squared distance between the two velocity fields on identical noised inputs,
\begin{equation}
  R_{\mathrm{vf}}(\theta) \;=\;
  \mathbb{E}_{s\sim\text{rollout},\,t,\,\varepsilon}
  \big\Vert\, v_\theta(x_t,t\mid s) - v_{\mathrm{SFT}}(x_t,t\mid s)\,\big\Vert^2 ,
  \label{eq:rvf}
\end{equation}
added to the loss with weight $\beta$. This is the tractable analogue of the per-denoising-step KL
that noise-space methods get for free, and it costs one frozen forward pass. Because the live and
frozen policies share normalization and see the same $(\varepsilon,t)$, the penalty is exactly $0$
at initialization and grows only as the policy drifts. The total per-iteration objective is
\begin{equation}
  \mathcal{L}(\theta) \;=\;
  \underbrace{\mathbb{E}_b\!\left[\, w_b \cdot \ell_{\mathrm{fm}}(x_b;\theta)\,\right]}_{\text{reward-weighted flow-matching BC}}
  \;+\; \beta\, R_{\mathrm{vf}}(\theta),
  \label{eq:loss}
\end{equation}
with $\ell_{\mathrm{fm}}$ the per-sample SmolVLA flow-matching MSE. We hold the anchor fixed at
$\pi_0$ across all iterations and use $\beta=10$ (a smoke sweep showed $\beta=50$ over-restrains,
and $\beta=0$ is the pure-filtering baseline). The implementation captures $v_\theta$ as the output
of \texttt{action\_out\_proj} via a forward hook, with no reimplementation of the forward pass.
There is also an AWR reading: the reference pull is partly implicit in fitting weighted samples
drawn from a policy that starts as $\psft$, and the explicit anchor holds it across iterations as
$\pi_k$ drifts.

\paragraph{The demo-seed remedy.}
Filtered BC needs some successes to bootstrap. A goal-only policy on \libten{} scores about 0\%
zero-shot (Section~\ref{sec:results}), so the filter starves. The remedy is to first SFT a
\libten{} specialist from \texttt{smolvla\_base}, rather than from the goal SFT, which avoids a
goal-transfer confound and lifts the policy off the floor so the loop has something to work with.

\paragraph{The collection-equals-evaluation protocol.}
The rollout success predicate equals LeRobot's $\pcsucc$ (we validate this as a smoke gate), so
each iteration's collection \emph{is} that policy's per-task evaluation. Seeding rollouts
deterministically (from 1000) makes the per-task numbers a paired evaluation: $\pi_0$, $\pi_1$, and
$\pi_2$ see identical initial states, so per-task deltas are not confounded by initialization. This
yields a paired $n{=}500$ evaluation on \libgoal{} (50 per task) and $n{=}20$ per task on \libten{}
at no extra cost.

\section{Experimental Setup}
\label{sec:setup}

All runs use lerobot 0.5.1, torch 2.10.0+cu128, transformers 5.3.0, python 3.12.13. Hardware is
matched to the bottleneck: an A100-80GB for the GPU-bound SFT and fine-tune steps, and an L4 for the
sim/CPU-bound rollouts and evaluation, where a larger GPU only idles through MuJoCo. Evaluation uses
relative control, \texttt{n\_action\_steps}=10, the camera \texttt{rename\_map}, and one process per
task; a multi-task eval crashes at the task transition with an EGL context-teardown error on
long-running runtimes. The eval CI at $n{=}100$ is about $\pm$7--10 points once within-task
clustering is accounted for, so we report per-task throughout and treat $n{=}500$ as the tighter,
paired measurement. Total project compute fits a \$20--60 budget on preemptible Colab Pro+, with
per-2000-step checkpointing to Drive for preemption resilience.

\section{Results}
\label{sec:results}

Table~\ref{tab:headline} collects the headline numbers; Figures~\ref{fig:traj}--\ref{fig:sft}
show the trajectories, the per-task structure, and the SFT contribution.

\begin{table}[t]
\centering
\caption{Headline results. All success rates are $\pcsucc$ (\%). RL trajectories are paired
(identical seeds across iterations). The \libten{} SFT figure (50.0, $n{=}100$) and its paired RL
baseline ($\pi_0=44.0$, $n{=}20$ per task, different seeds) measure the same checkpoint under
different protocols.}
\label{tab:headline}
\small
\begin{tabular}{@{}llll@{}}
\toprule
Stage & Suite & $n$ & Result \\
\midrule
Official baseline & \libgoal{} & 100 & 85.0 \\
Goal SFT (ours, 16k/20k) & \libgoal{} & 100 & 89.0 \quad($+4.0$, within noise) \\
RL (filtered BC, $\beta{=}10$) & \libgoal{} & 500 (paired) & $89.0 \to 91.0 \to 89.4$ \\
Zero-shot probe (goal-only) & \libten{} & 10/task & $\sim$0 (starvation) \\
Demo-seed SFT & \libten{} & 100 & 50.0 \\
RL (filtered BC, $\beta{=}10$) & \libten{} & 20/task (paired) & $44.0 \to 47.5 \to 42.5$ \\
\bottomrule
\end{tabular}
\end{table}

\begin{figure}[t]
\centering
\includegraphics[width=\linewidth]{figure1_trajectories}
\caption{Overall success across RL iterations. In both regimes the trajectory rises one iteration
and returns to (or below) the $\pi_0$ baseline (dashed) the next, so compounding is falsified.
\emph{Left:} saturated \libgoal{} (paired $n{=}500$). \emph{Right:} unsaturated, demo-seeded
\libten{} (paired $n{=}20$ per task).}
\label{fig:traj}
\end{figure}

\subsection{Baseline and SFT (validated pipeline)}
The official checkpoint scores 85.0\% on \libgoal{} ($n{=}100$), and no task scores 0\% (the floor
is 60\%), which rules out the control-mode and camera-key failure modes that would crater every task
uniformly. Our goal-only SFT reaches 89.0\% (Figure~\ref{fig:sft}, left). The $+4.0$ is
directionally consistent with goal specialization but sits inside the $n{=}100$ CI, so it is not a
claimed improvement. The SFT deliverables are a validated end-to-end pipeline and a clean checkpoint
$\psft$ to launch RL from.

\subsection{RL on saturated \libgoal{}: within-noise redistribution}
The paired $n{=}500$ trajectory is $89.0 \to 91.0 \to 89.4$ (Figure~\ref{fig:traj}, left).
Iteration 1 looks positive ($+2.0$), concentrated on the low-headroom tasks (task 9 $78\to90$, task
4 $88\to100$). Iteration 2 does not compound; it regresses $-1.6$ and erases those gains (task 9
$90\to80$, task 3 $86\to78$). $\pi_2$ is statistically indistinguishable from $\pi_0$ ($+0.4$). On a
near-saturated suite the method redistributes success across tasks within noise rather than
producing a durable aggregate gain (Figure~\ref{fig:pertask}, top).

\subsection{Zero-shot \libten{}: the starvation failure mode}
Evaluated zero-shot on \libten{}, the goal-only $\pi_0$ scores 0/10 on each probed task. With zero
successful rollouts there is nothing for filtered BC to imitate, so the loop cannot bootstrap. This
is the sparse-reward starvation failure, and it appears exactly where the headroom is largest. The
slowness corroborates it: every episode runs to the long-horizon timeout.

\subsection{Demo-seed remedy and RL on \libten{}}
A \libten{} specialist SFT lifts success from about 0\% to 50.0\% ($n{=}100$; Figure~\ref{fig:sft},
right), which removes the starvation. Running the same RL loop on this genuine headroom reproduces
the \libgoal{} pattern: paired $44.0 \to 47.5 \to 42.5$ (Figure~\ref{fig:traj}, right). Iteration
1's $+3.5$ is a marginal-task rescue, concentrated on tasks $\pi_0$ seeded with a few successes
(task 3 $60\to75$, task 7 $25\to40$, task 4 $10\to20$); iteration 2 reverses exactly those tasks
(task 3 to 55, task 7 to 25, task 4 to 10), landing $-1.5$ below $\pi_0$. Establishing the same null
on an unsaturated, bootstrappable suite makes it stronger and more general than \libgoal{} alone.

\begin{figure}[t]
\centering
\includegraphics[width=\linewidth]{figure2_pertask}
\caption{Per-task success for $\pi_0/\pi_1/\pi_2$. \emph{Top:} \libgoal{} ($n{=}500$); gains and
losses cancel across tasks. \emph{Bottom:} \libten{} ($n{=}20$ per task); the marginal tasks
rescued in iteration 1 (t3, t4, t7) reverse in iteration 2, while the fully-starved task 8 (0 seed
successes) stays on the floor.}
\label{fig:pertask}
\end{figure}

\section{Discussion}
\label{sec:discussion}

\paragraph{Two mechanisms.}
The per-task data show why the aggregate does not move durably. First, rescue-then-reversal:
lightweight RL transiently lifts weak tasks that the SFT policy seeded with at least a few
successes, because the filter then has something to amplify, but the next iteration collects from an
already-drifted policy under fixed seeds and reverses them. The up-then-down per-task signature is
regression to the mean, or stochastic redistribution, not a learned gain. Second, the starvation
floor: a task with zero seed successes (\libten{} task 8) cannot be bootstrapped by a filter at all
and stays at 0 across iterations. The lone iteration-2 ``10'' is 2/20 seed luck rather than a
bootstrap, since $\pi_1$ had no task-8 successes to imitate.

\paragraph{Why the null is informative.}
A negative result is most useful when an obvious confound does not explain it away. This one
survives two natural objections. It is not just a ceiling effect, because it also holds on
unsaturated \libten{} with real headroom, and it is not just starvation, because the demo-seed
removes starvation and the gain still does not compound. What lightweight RL produces is a
redistribution of success across tasks within measurement noise, in both regimes tested.

\paragraph{Limitations.}
The main ones. (i) Single seed per policy: the per-iteration swings are about 1 SD, so it is the
up-then-down trajectory, not any single delta, that reads as noise. (ii) The \libten{} measurement
is $n{=}20$ per task (binomial CI about $\pm$20), so the reliable figure is the aggregate. (iii) A
single anchor strength $\beta{=}10$ rather than a full sweep. (iv) RL starts from a 16k
(under-trained) SFT base, but the SFT-to-RL delta is measured against that same fixed base on both
sides, so it does not contaminate the RL delta. (v) Sim-only, outcome-only binary reward; under it,
``reward-weighted'' BC equals filtered BC plus a mild time-to-success weight. (vi) Few outer
iterations. A claim-grade result would need $n{=}500$ with at least three seeds varying both
training and evaluation seeds, tens of hours on \libten{}'s slow renderer, which we scope out and
flag rather than approximate.

\begin{figure}[t]
\centering
\includegraphics[width=\linewidth]{figure3_sft}
\caption{The SFT contribution. \emph{Left:} goal-only SFT vs.\ the official baseline on \libgoal{}
($+4.0$, within noise). \emph{Right:} the demo-seed SFT lifts \libten{} off the $\sim$0\% goal-only
zero-shot floor to 50.0\%, which enables, but does not guarantee, an RL effect.}
\label{fig:sft}
\end{figure}

\section{Conclusion and Future Work}

We asked whether a tractable, flow-matching-compatible RL step (reward-weighted filtered BC with a
vector-field trust-region anchor) improves a SmolVLA VLA beyond SFT on LIBERO, and tested it as
carefully as a \$20--60 budget allows. Across a saturated in-distribution suite and an unsaturated
demo-seeded one, the answer is no: success redistributes within noise rather than improving durably,
through a rescue-then-reversal mechanism and a starvation floor. The deliverables are the method,
the paired collection-equals-evaluation protocol, and a map of where this class of method can and
cannot move a flow-matching policy. Natural next steps are a claim-grade evaluation ($n{=}500$, at
least three seeds) to put a confidence interval on the flat delta; an anchor-strength ($\beta$)
sweep; flow-matching-native on-policy methods that recover a tractable
likelihood~\citep{ren2024dppo,chen2025pirl,reinflow2025,flowgrpo2025}, which target generalization
rather than in-distribution success; and real-robot validation. The null, with its mechanisms, is
the result.

\paragraph{Reproducibility.} All six result files (\texttt{results/*.json}) are released with the
exact per-task numbers, together with the figure-generation script, the three RL scripts (rollout,
fine-tune, eval), and the SFT runbook with the five fixes. The deterministic list of trained
episode indices accompanies the checkpoints.

\small
\bibliographystyle{plainnat}
\bibliography{references}

\end{document}
